\chapter{Introduction}

Coastal regions are home to a significant portion of the global population,
and they are increasingly affected by human activities. Desalination processes \cite{omerspahic2022brine} have
become the main solution to provide freshwater to many of these regions, and even though
more efficient techniques have been developed, the accompanying environmental issues have 
been overlooked, in particular the consequences of the discharge of desalination 
by-products, such as the effects of concentrate effluent on the marine coastal ecosystem. 
Moreover, other industrial activities, such as offshore oil and gas operations \cite{cordes2016oilgas}, 
have been shown to have a significant impact on the marine environment, including the
release of pollutants and the disruption of marine habitats. 
Therefore, the preservation and monitoring of marine environments is a strategic priority for coastal regions, 
especially in the Arabian Gulf, where the combination of high salinity, warm temperatures, and human activities 
poses a significant threat to marine ecosystems \cite{sheppard2010gulf}.

Coral reefs provide an existential habitat for a wide range of marine species,
offering shelter, food and breeding grounds. However, they are highly sensitive
to environmental changes, and their degradation could have severe consequences for the 
environment and the economy. Coral restoration is advocated as a promising tool to reduce coral
loss and restore damaged reefs. Many restoration techniques have been developed and many 
studies have been conducted to evaluate their effectiveness, but the success of these methods
is still limited \cite{mula2025restoration}. $30-40\%$ of projects fail due to poor planning, 
unrealistic objectives, inadequate regular maintenance and persistent anthropogenic pressures. 
Moreover, restoration sites are often clustered around human settlements, which are not always 
the most suitable locations for coral growth, but limits human hazards and encourages safety 
for divers, while also reducing costs. Finally, the search areas for suitable restoration sites
are enormous, limiting the possibility search and monitoring due to low availability of divers,
autonomy and time.

For these reasons, the use of autonomous underwater vehicles (AUVs) has been proposed as a solution
to overcome the limitations of human divers and to increase safety and efficiency in the search task. 
Moreover, given the unknown and complex nature of the marine environment, AUVs need to be equipped with
advanced sensors and algorithms to navigate and map the environment. To this end, Multi-agent 
Reinforcement Learning (MARL) algorithms have been recently proposed and increasingly used to overcome
this task, enhancing the capabilities of Reinforcement Learning (RL) algorithms with swarm and cooperation 
capabilities, so that multiple agents can work together to achieve a common goal. A previous thesis \cite{insaghi2025thesis}
developed and analyzed a PPO-PSO algorithm to efficiently use swarm capabilities to overcome 
local optima and improve the performance of the PPO algorithm. The results showed that the PPO-PSO 
algorithm outperformed the standard PPO algorithm in terms of convergence speed and final performance, 
demonstrating the potential of MARL algorithms for complex tasks in unknown environments. However,
the environment used in the previous thesis was a simplified simulation, and the necessity to test 
the algorithm in a more realistic environment is the main motivation for this thesis. Therefore,
the main objective of this thesis is to evaluate the performance of a suitable MARL algorithm in a 
more realistic environment, where the agents are required to navigate and search for suitable 
coral restoration sites in a complex and unknown environment.

This thesis is organized as follows: Chapter 2 provides a comprehensive theoretical background on the 
Reinforcement Learning algorithms, including the multi-agent extension and general problem formulation. Chapter 3 
describes the problem more specifically, including the assumptions, constraints and tools used. Chapter 4 
introduces the proposed MARL algorithm, including its architecture, hyperparameters and implementation details.
Chapter 5 presents the experimental setup, where multiple scenarios are created to test the capabilities 
of the obtained policy and its generalization capabilities. Finally, Chapter 6 concludes the thesis,
summarizing the main findings, discussing the limitations and proposing future research directions.